\section{Capítulo~\ref{sec:neurontypes}. Tipos de neurônios}

Os números do capítulo se apoiam em contagens diretas de células no cérebro humano, onde elas existem; a regra ``um neurônio, um neurotransmissor'', em atlas celulares e em experimentos que também encontraram as exceções; o papel de \nt{DOPA}, em registros de disparos; e os papéis dos demais canais de difusão, em hipóteses.

{\footnotesize
\setlength{\tabcolsep}{3pt}\renewcommand{\arraystretch}{1.15}
\begin{longtable}{>{\raggedright}p{0.5cm}>{\raggedright}p{4.0cm}>{\raggedright}p{5.6cm}>{\raggedright}p{2.3cm}>{\raggedright\arraybackslash}p{1.7cm}}
\toprule
N.º & Proposição & Experimento e o que foi medido & Fonte & Status \\
\midrule\endfirsthead
\toprule
N.º & Proposição & Experimento e o que foi medido & Fonte & Status \\
\midrule\endhead
1 & O cérebro humano tem cerca de $8.6\cdot10^{10}$ neurônios, e quase todos os neurônios \nt{GLUT} são pequenos neurônios do cerebelo (tab.~\ref{tab:types}) & Fracionador isotrópico: o tecido é dissolvido numa suspensão homogênea de núcleos, e contam-se os núcleos com o marcador neuronal NeuN. Em homens adultos, $86.1\pm8.1$ bilhões de neurônios; apenas $19\,\%$ deles estão no córtex cerebral, a maior parte está no cerebelo & \cite{Azevedo2009} & medido \\
2 & Quase todo neurônio tem um neurotransmissor principal (proposição~\ref{prop:dale}), mas há exceções & Atlas transcriptômico do cérebro do camundongo: cerca de $4\cdot10^6$ células, $5322$ clusters; a cada cluster se atribui um neurotransmissor pelos genes de síntese e empacotamento. As exceções foram medidas diretamente: neurônios \nt{DOPA} em fatias liberam também GABA pelo mesmo transportador vesicular e inibem neurônios do estriado; em parte dos axônios \nt{DOPA}, glutamato e dopamina saem de terminais diferentes do mesmo fio (microscopia eletrônica e optogenética) & \cite{Strata1999,Yao2023,Tritsch2012,Zhang2015} & medido \\
3 & Toda a coluna da matriz de pesos tem o mesmo sinal~\eqref{eq:dalesign} & Dedução no capítulo a partir da proposição~\ref{prop:dale} e do fato de que os receptores de glutamato são quase todos excitatórios, e os de GABA, inibitórios. Não há verificação direta de ``todos os destinatários de um neurônio recebem peso do mesmo sinal'' como regra geral; contraexemplos: a liberação conjunta de GABA por neurônios \nt{DOPA} (linha~2) e destinatários com receptores de sinais diferentes (linha~11) & dedução no capítulo & teoria \\
4 & De três quartos a quatro quintos dos neurônios do córtex são \nt{GLUT}, de um quinto a um quarto são \nt{GABA} & Imunomarcação para GABA em colunas de $50$\,\textmu m de largura por toda a espessura do córtex em $10$ áreas corticais de macaco: em $8$ áreas os neurônios GABAérgicos são cerca de $25\,\%$ de todos os neurônios; no córtex visual primário, menos de $20\,\%$ & \cite{Hendry1987} & medido \\
5 & Um neurônio \nt{GLUT} tem cerca de $10^4$ saídas & Estereologia em autópsias: no neocórtex de homens jovens, $1.64\cdot10^{14}$ sinapses (microscopia eletrônica, dissector) e cerca de $2.3\cdot10^{10}$ neurônios. A razão dá $\sim7\cdot10^3$ sinapses por neurônio; em média, o número de entradas é igual ao de saídas & \cite{Tang2001synapses,Pakkenberg1997} & medido \\
6 & A saída de um neurônio \nt{DOPA} se ramifica em $10^5$--$10^6$ terminais; é uma estimativa pela cobertura do alvo, não uma contagem & Marcação viral da membrana de neurônios \nt{DOPA} isolados de rato e reconstrução completa do axônio: um neurônio cobre de $0.45$ a $5.7\,\%$ (em média $2.7\,\%$) do volume do estriado. Mediu-se a cobertura, não o número de terminais: este foi deduzido da cobertura em ordem de grandeza, sem contagem direta dos terminais. Para \nt{SERO} e \nt{NORA}, os números de terminais são estimativas grosseiras & \cite{Matsuda2009} & medido \\
7 & Os neurônios de difusão são poucos: \nt{DOPA} $\sim5\cdot10^5$ (o principal dos dois grupos, limite inferior), \nt{SERO} $\sim3\cdot10^5$ (soma de duas contagens parciais), \nt{HIST} $\sim6\cdot10^4$ (tab.~\ref{tab:types}, fig.~\ref{fig:typepie}) & Contagens estereológicas em humanos: $550\,000$ neurônios pigmentados na substância negra (sem a área tegmental ventral); $235\,000$ neurônios no núcleo dorsal da rafe (todos os neurônios do núcleo) e mais $125\,000$ neurônios serotoninérgicos no tegmento pontino; cerca de $64\,000$ neurônios histaminérgicos do hipotálamo. Para \nt{NORA}, \nt{GLYC}, \nt{ACET} e \nt{ADRE} não há contagem direta: na tabela são estimativas grosseiras de ordem de grandeza & \cite{Pakkenberg1991,Baker1990,Baker1991,Airaksinen1991} & medido \\
8 & O glutamato na fenda sobe até cerca de um milimolar e cai em $\sim1\ms$ & Pelo deslocamento de um bloqueador competitivo de dissociação rápida dos receptores NMDA durante a transmissão sináptica (cultura de neurônios do hipocampo): pico de $1.1$\,mM, constante de decaimento de $1.2\ms$ & \cite{Clements1992} & medido \\
9 & No córtex em funcionamento, as correntes excitatória e inibitória estão quase equilibradas, e o neurônio fica perto do limiar & Teoria: uma rede com conexões fortes chega sozinha ao regime equilibrado. Experimento: no córtex pré-frontal do furão in vivo, durante os estados ``up'', o potencial de reversão da corrente sináptica se mantém em torno de $-37$\,mV por centenas de milissegundos, embora a condutância varie em $\sim21$\,nS: excitação e inibição sobem e descem juntas & \cite{vanVreeswijk1996,Haider2006} & medido \\
10 & Os neurônios \nt{DOPA} informam o erro de predição de recompensa~\eqref{eq:rpe} (fig.~\ref{fig:rpe}) & Disparos de neurônios \nt{DOPA} de macaco: rajada ao suco inesperado; após o aprendizado, rajada ao sinal e nenhuma resposta ao suco previsto; queda da frequência quando o suco prometido não é dado. Num experimento quantitativo, a frequência é proporcional à diferença entre a recompensa atual e a média ponderada das anteriores, mas só para erros positivos. A equação~\eqref{eq:rpe} em si é o algoritmo de diferenças temporais & \cite{Schultz1997,Bayer2005,SuttonBarto2018} & medido \\
11 & O sinal e a velocidade são definidos pelo receptor: ionotrópicos, frações de milissegundo; metabotrópicos, muito mais lentos (proposição~\ref{prop:receptor}) & Pulsos rápidos de glutamato em fragmentos de membrana de neurônios do hipocampo: a corrente pelos receptores ionotrópicos sobe (de $20$ a $80\,\%$) em $0.2$--$0.6\ms$ e decai em $2$--$3\ms$. O potencial inibitório lento das células piramidais é abolido por um bloqueador seletivo do receptor metabotrópico de GABA. Duas famílias de receptores de dopamina com ação oposta; no estriado, os neurônios com receptor D1 precisam de dopamina para reforçar sinapses, e os com D2, para enfraquecê-las & \cite{Colquhoun1992,DutarNicoll1988,Beaulieu2011,Shen2008} & medido \\
12 & O canal de difusão não é um fio só, mas um feixe de poucas linhas (seção~\ref{sec:buses}) & Marcação viral e genética dos neurônios serotoninérgicos do núcleo dorsal da rafe do camundongo: os neurônios que enviam saída para a amígdala e para o córtex frontal ficam em lugares diferentes do núcleo, ramificam-se em regiões complementares e respondem de modo oposto a um estímulo aversivo. Revisão: subgrupos de neurônios do locus coeruleus (\nt{NORA}) codificam processos diferentes & \cite{Ren2018,Poe2020} & medido \\
13 & \nt{NORA} define o ganho $g$ & Hipótese do ganho adaptativo; modelo de rede em que as catecolaminas aumentam a inclinação da característica de transferência. Não há medição direta de ``$g$ cresce com a noradrenalina'' na rede inteira; mediu-se que o diâmetro da pupila acompanha os disparos dos neurônios do locus coeruleus do macaco & \cite{AstonJones2005,ServanSchreiber1990,Joshi2016} & hipótese \\
14 & \nt{ACET} alterna ``escrita/leitura'' e define a taxa de aprendizado & Hipótese. Base: em fatias de córtex olfatório de rato, agonistas da acetilcolina ($100$\,\textmu M) enfraquecem as sinapses internas, de ``lembrança'', em mais de $60\,\%$, e as de entrada em menos de $15\,\%$ & \cite{Hasselmo2006,HasselmoBower1992} & hipótese \\
15 & \nt{SERO} define o fator de desconto $\gamma$; os neurônios serotoninérgicos trabalham de forma regular, alguns disparos por segundo, e quase silenciam numa das fases do sono & Hipótese ``serotonina $=\gamma$''. O argumento mais forte: a ativação optogenética dos neurônios serotoninérgicos do núcleo dorsal da rafe do camundongo reduz o número de desistências prematuras na espera de uma recompensa adiada. A frequência de disparos e o silêncio no sono com movimentos oculares rápidos foram medidos (revisão de registros) & \cite{Doya2002,Miyazaki2014,JacobsAzmitia1992} & hipótese \\
\bottomrule
\end{longtable}}
